\section{Related work}\label{sec:related}

\subsection{Graph collaborative filtering}

LightGCN \citep{he2020lightgcn} is the backbone used here. It is one of several graph
convolutional recommenders, which include NGCF \citep{wang2019ngcf}. Other graph models
include UltraGCN \citep{mao2021ultragcn}, and self-supervised and contrastive graph variants
include SGL \citep{wu2021sgl} and SimGCL \citep{yu2022simgcl}. Linear and closed-form models
include EASE \citep{steck2019ease}, GF-CF \citep{shen2021gfcf} and BSPM \citep{choi2023bspm}.
Implicit matrix factorisation with weighted least squares, iALS \citep{hu2008ials,rendle2022ials},
is a standard baseline, and carefully tuned simple methods often match published graph
models \citep{anelli2023myth}. A survey of graph-based recommenders is given by
\citet{wu2022gnnsurvey}. These models fit or propagate over the logged graph and do not
model which pairs were exposed. The present method keeps the propagation of LightGCN and
changes the training signal.

\subsection{Debiasing under missing-not-at-random feedback}

Ratings and implicit feedback are missing not at random
\citep{marlin2009collaborative}, and surveys organise the resulting biases and the
causal methods that address them \citep{chen2023bias,gao2024causalsurvey}. Propensity
methods re-weight the loss by inverse propensities or self-normalised estimators
\citep{schnabel2016,swaminathan2015snips,saito2020unbiased}. Steck \citep{steck2010mnar}
showed that missing ratings carry information, proposed test measures that can be
estimated without bias when ratings are missing not at random, and derived
training objectives that account for the missing entries. Doubly robust learning adds an imputation model
\citep{wang2019drjl}, and later variants reduce its bias or variance
\citep{guo2021mrdr,dai2022gdr,li2023stabledr}. Information-theoretic counterfactual
learning from MNAR feedback is studied by \citet{wang2020cvib}. Some debiasing methods use a small uniform
sample for training, as in AutoDebias \citep{chen2021autodebias}, CausE
\citep{bonner2018cause} and KDCRec \citep{liu2020kdcrec}, which distils from a model fitted on uniform data into one trained on the biased log.
Work on popularity bias includes model-agnostic counterfactual reasoning \citep{wei2021macr},
causal intervention \citep{zhang2021pda}, disentangled user interest and conformity
\citep{zheng2021dice}, popularity debiasing that connects user and item perspectives
\citep{boratto2021connecting}, flexible manipulation of popularity bias
\citep{okamura2024flexibly} and invariant debiasing learning via biased imputation
\citep{bai2025invariant}. The present method applies propensity and DR estimation to the
nuisances that define which pairs act as negatives.

\subsection{False negatives and negative weighting in implicit feedback}

This is the closest group of methods. ExpoMF \citep{liang2016exposure} treats exposure
as a latent variable and infers an exposure posterior for unclicked items within a
matrix factorisation model. Rel-MF \citep{saito2020unbiased} re-weights observations
by propensity, and unbiased pairwise learning \citep{saito2020ictir} applies the same
idea to pairwise objectives. Sampling-based methods treat false negatives directly.
SRNS \citep{ding2020srns} identifies likely false negatives in the sampling step, and
debiased contrastive learning \citep{chuang2020dcl} corrects the sampling of negatives
that are in fact positives. In contrastive collaborative filtering, the BC loss
\citep{zhang2022bc} adds a margin that depends on the popularity-bias degree of each
interaction, and AdvInfoNCE \citep{zhang2023advinfonce} assigns each negative an
adversarially learned hardness weight. Correct-and-Weight \citep{yin2026cw} combines a loss
correction with a weighting of negatives against false negatives, and is tested with MF,
LightGCN and XSimGCL on random splits. Exposure-aware negative sampling for graph
recommenders uses in-community item popularity as a proxy for exposure \citep{chen2026icpns}.
The present weight $1-Q$ has the posterior form of ExpoMF but is computed in closed form
from the nuisances of a DR estimator. It is applied to a graph model with a pairwise
loss and is combined with a prior offset. The original formulations of ExpoMF, SRNS and
AdvInfoNCE do not use a propensity and an imputation of this kind.

\subsection{Debiasing inside graph models}

NAVIP \citep{kim2022navip} weights neighbour aggregation by inverse propensities inside
a graph neural recommender. IPS-weighted BPR on LightGCN has been studied on simulated
data \citep{raja2025ipsbpr}. r-AdjNorm \citep{zhao2022radjnorm} tunes the normalisation
of the adjacency matrix to control popularity amplification, and how graph convolutions
amplify popularity bias is the subject of \citet{chen2024gcnpop}. APDA
\citep{zhou2023apda}, CAGED \citep{que2025caged} and DPAA \citep{dpaa2026} learn
aggregation weights or debiased message passing. AdvDrop \citep{zhang2024advdrop} learns an
adversarial edge dropout on the interaction graph and reports IPS and DR baselines on
Coat, Yahoo and KuaiRec. PPD \citep{islam2026ppd} removes a popularity direction from
pre-trained LightGCN embeddings without retraining, and is evaluated on KuaiRec, Coat and
Yahoo! R3 with unbiased test data. DR-GNN \citep{wang2024drgnn} applies distributionally
robust optimisation over graph neighbourhoods; its exposure-shift experiments use Coat and
Yahoo, and it does not use propensities or imputations. These methods change the
aggregation, the graph, the embeddings or the robust training objective. The present method
leaves the propagation of LightGCN unchanged. It changes which pairs count as negatives in
the loss and adds a prior to the score.

\subsection{Score offsets and priors}

Adding a fixed offset to a score is a familiar device. In long-tail classification,
logit adjustment adds a class prior to the logits \citep{menon2021logit}, and in
generalised linear models an offset is a known term of the linear predictor. The offset
used here is different in one respect: it is a debiased estimate of population-level
preference, obtained from an inverse-propensity-weighted imputation, and not a
class prior.

\begin{table}[p]
\centering
\caption{Positioning of the method relative to the main related families. Entries are
qualitative. ``Acts on'' names the part of the model that the correction changes.
``Uniform data'' indicates whether a uniformly collected sample enters training; where
a method uses it only for validation, the cell says so. For
the methods that we did not run (APDA, CAGED, AdvDrop, PPD, DPAA, DR-GNN, Correct-and-Weight) the entries follow their published
descriptions. Coat ships its propensities, which were estimated from a uniform sample
by the dataset authors; this affects the propensity-based rows on Coat only.}
\label{tab:positioning}
\scriptsize
\adjustbox{max width=\textwidth}{%
\begin{tabular}{p{3.2cm}p{1.9cm}p{1.8cm}p{1.6cm}p{2.6cm}p{1.8cm}}
\toprule
Method family & Graph backbone & Propensity & Imputation & Acts on & Uniform data \\
\midrule
LightGCN, SimGCL & yes & no & no & none (BPR or contrastive loss) & no \\
NAVIP & yes & yes & no & adjacency & no \\
r-AdjNorm, APDA, CAGED, DPAA & yes & no & no & adjacency or aggregation & no \\
AdvDrop & yes & no (IPS-CN baseline only) & no (DR baseline only) & graph (edge dropout) & no (not stated) \\
PPD & yes & no & no & embeddings (post hoc) & no (validation only) \\
DR-GNN & yes & no & no & robust objective over neighbourhoods & not stated \\
Correct-and-Weight & yes & no & no & loss (correction and negative weights) & no \\
IPS, SNIPS, DR, DR-JL, MRDR & no (matrix factorisation) & yes & DR only & loss & no (as described) \\
ExpoMF, Rel-MF, UBPR & no & latent or yes & no & loss or exposure posterior & no \\
AdvInfoNCE, BC loss & varies & no & no & loss (negative weights or margins) & no \\
AutoDebias & no (matrix factorisation) & learned weights & learned & loss (weights and imputation) & yes \\
CausE & no (matrix factorisation) & no & no & embeddings (discrepancy regulariser) & yes \\
DRUP (walk correction) & yes & yes & yes & multi-hop adjacency & no \\
DR adjacency (no correction) & yes & yes & yes & adjacency & no \\
\midrule
\qp{} (this paper) & yes & yes & yes & negatives and score & no (Coat: shipped propensities) \\
\bottomrule
\end{tabular}}
\end{table}

\subsection{Evaluation protocols and the gap}

Progress claims in top-$N$ recommendation depend on the strength of the baselines and on
the evaluation protocol \citep{dacrema2019progress,rendle2019baselines,anelli2023myth}.
Inference for the generalisation error under resampling is treated by \citet{nadeau2003}.
Sampled
evaluation metrics can also be inconsistent with their exact counterparts
\citep{krichene2020sampled}. Section~\ref{sec:setup} describes the protocol used here.

The uses of uniform data differ across debiasing papers. CDR \citep{song2023cdr} uses 10\%
of the unbiased data for validation and 90\% for evaluation on Coat, Yahoo! R3 and KuaiRand,
with a matrix-factorisation backbone. DCE \citep{kweon2024dce} holds out 10\% of the biased
training set for validation, does not use the unbiased test set for training, and evaluates on
Coat, Yahoo! R3 and KuaiRec. BISER \citep{lee2022biser} trains on MNAR data and evaluates on
MAR data on Coat and Yahoo! R3, with 30\% of each user's training ratings held out for
validation. ConAlign \citep{zhang2026conalign}, an arXiv preprint, trains a dual-tower model in
which an unbiased tower is trained on uniform data, on Coat, Yahoo! R3 and KuaiRand-Pure.

Three lines of work meet in this paper but have rarely been combined. Doubly robust and
propensity methods (IPS, DR-JL, MRDR) are pointwise and are usually applied to matrix
factorisation. Graph recommenders that address exposure or popularity bias (NAVIP,
r-AdjNorm, APDA, CAGED, AdvDrop, PPD, DPAA) act on the adjacency, the aggregation, the graph
or the embeddings rather than on the sampled negatives. Exposure-aware and false-negative-aware
methods (ExpoMF, SRNS, DCL, AdvInfoNCE) do not use the propensity and imputation that DR
estimation provides. Partial overlaps also include IPS-weighted BPR on LightGCN with simulated
propensities \citep{raja2025ipsbpr} and exposure-aware negative sampling with a popularity proxy
\citep{chen2026icpns}; neither uses an imputation or an unbiased-test evaluation. We are not aware of a previous study that uses a propensity and an
IPS-fitted imputation to build a posterior for non-logged pairs and uses it to weight the
sampled negatives of a graph recommender's pairwise loss, or that adds a debiased imputation
prior to a graph model's score. We are also not aware of a study that evaluates a graph model
on KuaiRand-Pure together with Coat, Yahoo! R3 and KuaiRec, with uniform data used only for
validation and testing. Coat, Yahoo! R3 and KuaiRec are already used for graph debiasing by
AdvDrop, PPD and DPAA; we found no graph study that uses KuaiRand. The novelty claimed here is
therefore limited to the Q-posterior weighting of negatives and the imputation-logit offset,
both inside the pairwise loss of a graph model. We did not find a study that combines these two
in that way. The search used general web search engines, not bibliographic databases, so the
absence of a hit is not proof of absence.

Several studies report that tuned simple baselines are competitive with published recommenders
\citep{dacrema2019progress,rendle2019baselines,anelli2023myth}.
